\section{Introduction}\label{sec:intro}

Many real time series are not sampled on a regular grid. Sensors fail or
skip readings, medical measurements are taken when a patient needs them,
and user activity is recorded only when something happens. The usual
neural models for sequences, the LSTM \cite{hochreiter1997lstm} and the
GRU \cite{cho2014gru}, read one value per step and have no notion of the
time between two values, so irregular data has to be filled in first or
the time gap has to be added as an extra input.

Two families of methods use the observation times directly. The first
consists of neural networks whose hidden state evolves in continuous
time: liquid time-constant networks (LTC) \cite{hasani2021ltc}, their
closed-form approximation (CfC) \cite{hasani2022cfc}, and neural controlled
differential equations (Neural CDE) \cite{kidger2020ncde}. The second are
mathematical features of the series followed by a simple classifier: path
signatures \cite{chen1958integration, lyons2022sigml}, which describe the
shape of the path and accept uneven times, and dynamic mode decomposition
\cite{schmid2010dmd}, which summarises the series by the eigenvalues of a
linear model of its dynamics. The features are much cheaper to compute
than the networks are to train. This work compares the two families under
one protocol and asks three questions:
\begin{itemize}
\item[RQ1] How do CfC, LTC and the Neural CDE compare to the LSTM and GRU on
regularly sampled time series?
\item[RQ2] What happens to each model when points are missing and the time
steps become uneven?
\item[RQ3] Can signature or DMD features with a simple classifier get close
to these networks, and how much cheaper are they?
\end{itemize}

To answer them, we use four public datasets of the UEA archive
\cite{bagnall2018uea} and remove 30, 50 or 70\% of the points of every
series at random, following the protocol of Kidger et al.\
\cite{kidger2020ncde}, with the same gaps for every method. All methods
are measured by accuracy, macro-F1, training time and model size. This
report describes the setup and the first results: the full experiment
for the feature methods and a first, untuned run of the neural models.
Section~\ref{sec:related} reviews related work and the published results
on these datasets, Section~\ref{sec:model} gives the mathematical model of
the data, the methods and the experiments, and Section~\ref{sec:case}
presents the case study.
